%% kapa: exam=1 round=19 year=2008 period=1 form=A subject=appraisal_law subject_ko=부동산관계법규 seq=38 number=78
%% source: Q-Net 감정평가사 기출문제 (cst00309) artlSeq=1001962 file=extracted/1008/1001962/2008년 제19회 감정평가사 1차시험 기출문제/1교시.hwp
%% conversion: kordoc 4.24.51 -> tools/gen1.py (auto); review=proofread-sonnet; flags=none
\begin{question}{78}
부동산등기법상 가등기에 관한 설명으로 옳지 않은 것은? (다툼이 있는 경우에는 판례에 의함)
\begin{choices}
  \item 부동산권리변동을 목적으로 하는 청구권이 시기부 또는 정지조건부일 경우뿐만 아니라 기타 장래에 확정될 것인 때에도 가등기를 할 수 있다.
  \item 부동산등기법상의 가등기는 물권 또는 부동산임차권의 변동을 목적으로 하는 청구권을 보전하기 위해서만 가능하므로 물권적 청구권을 보존하기 위해서는 할 수 없다.
  \item 가등기에 의해 보존된 청구권을 양도한 경우에 가등기에 대한 부기등기 형식으로 경료할 수 있다.
  \item 가등기에 의해 후일 본등기가 경료된 때에는 본등기에 의한 물권변동의 효력이 가등기한 때로 소급하여 발생한다.
  \item 가등기는 신청서에 가등기의무자의 승낙서 또는 가처분명령의 정본을 첨부하여 가등기권리자가 신청할 수 있다.
\end{choices}
\end{question}
\answer{4}
